% ===========================================================================
% WAVVY REPORT TEMPLATE - v.1.0.0
% ===========================================================================

\documentclass[a4paper, 11pt, table, dvipsnames]{article}

% ===========================================================================
% PACKAGES
% ===========================================================================

% Fonts and symbols
\usepackage{lmodern}
\usepackage{amsmath, amssymb}
\usepackage[utf8]{inputenc}
\usepackage[T1]{fontenc}
\usepackage{inconsolata}
\usepackage{wasysym}
\usepackage{microtype}

% Layout, page setup
\usepackage[margin=2.5cm, a4paper]{geometry}
\usepackage{titlesec}
\usepackage{appendix}
\usepackage{multicol}
\usepackage{multirow}
\usepackage{parskip}
\usepackage{tabto}
\usepackage{pdflscape}
\usepackage{enumitem}
\usepackage{adjustbox}
\usepackage[super]{nth}

% Science and language typesetting
\usepackage{bm}
\usepackage[
	arc-separator=\,,
	retain-explicit-plus,
	inter-unit-product=\cdot,
	detect-weight=true,
	detect-family=true,
	range-phrase=--,
	range-units=single
]{siunitx}
\usepackage[version=4]{mhchem}
\usepackage{nicematrix}
\usepackage[makeroom]{cancel}
\usepackage{circledsteps}
\usepackage[italian]{babel}
\usepackage{csquotes}

% Graphics, captions, tables
\usepackage{graphicx}
\usepackage{float}
\usepackage{tikz, tikz-3dplot}
\usepackage{pgfplots}
\usepackage{pdfpages}
\usepackage[
	justification=centering,
	labelfont={small,bf},
	font={small}
]{caption}
\usepackage{subcaption}
\usepackage{array}
\usepackage{tabularx}
\usepackage{booktabs}
\usepackage[outline]{contour}

% TikZ options setup
\usetikzlibrary{
	calc,
	arrows,
	arrows.meta,
	positioning,
	decorations.pathreplacing,
	decorations.markings,
	decorations.text,
	calligraphy,
	pgfplots.dateplot,
	shapes.geometric
}

% References and links
\usepackage{hyperref}
\usepackage[noabbrev]{cleveref}
\usepackage[nottoc, numbib]{tocbibind}

% Miscellaneous
\usepackage{datetime2}
\usepackage{iftex, etoolbox, expl3, xparse}

\pgfplotsset{compat=1.18}
\AtBeginDocument{\RenewCommandCopy\qty\SI}

% ===========================================================================
% METADATA
% ===========================================================================

\title{\textbf{Archetipo del brand}}
\author{\textbf{Andrea Zorzi} -- Wavvy}
\date{Versione 1.0.0 \\ \today}

\begin{document}
\maketitle

\tableofcontents

\newpage

% ===========================================================================
% BODY
% ===========================================================================

\section{Scopo e perimetro}

Questo documento descrive l'infrastruttura psicologia e l'archetipo
di brand dell'applicazione. Serve a definire le linee
guida strategiche inconfutabili per le future scelte di UI/UX
design, marketing e comunicazione, al fine di garantire
coerenza e riconoscibilità del brand.

\begin{itemize}
    \item Piattaforme coinvolte: Applicazione Mobile
\end{itemize}

\subsection{Perimetro}
Questo documento si concentra sui seguenti aspetti:
\begin{itemize}
    \item Identificazione e giustificazione dell'architettura
          archetipica
\end{itemize}

Sono esplicitamente esclusi:
\begin{itemize}
    \item La generazione di codici colore definitivi e la UI
          palette

    \item La definizione di linee guida di comunicazione e tone
          of voice

    \item L'implementazione dell'app

    \item La definizione di linee guida per il marketing e la
          comunicazione
\end{itemize}

\section{Panoramica architetturale}
L'architettura dell'archetipo viene definita in modo ibrido,
con un blend strategico tra più archetipi, per massimizzare
l'acquisizione utente e la ritenzione emotiva. Sulla base
dell'analisi di mercato (DACH Region), l'identità è strutturata
tra due modelli psicologici junghiani:
\begin{itemize}
    \item The Everyman (L'Uomo Comune)

    \item The Lover (L'Amante)
\end{itemize}
L'app si presenta al mondo come il grande unificatore
democratico, ma trattiene nel suo nucleo l'energia sensoriale,
l'intimità e la privacy.

\section{Principi di progettazione}
Wavy segue i seguenti principi di progettazione:

\subsection{Inclusione ed assenza di attrito}
L'app democratizza l'approccio eliminando il timore del
rifiuto faccia a faccia.

\paragraph{Direttiva UI/UX}
L'intefaccia utente non deve creare elitarismo. A
differenza delle piattaforme tradizionali, non devono esserci
metriche di popolarità aggressive (es. contatori di follower
in primo piano) o \textit{paywall} visibili che dividono gli
utenti in classi.

\paragraph{Obiettivo}
L'esperienza d'uso deve comunicare un senso di uguaglianza
e appartenenza al medesimo spazio fisico, rendendo l'interazione
naturale e priva di pressione sociale.

\subsection{Validazione della presenza reale}
L'app è progettata come infrastruttura digitale al servizio
del mondo tangibile.

\paragraph{Obiettivo}
Convertire l'incontro sull'app in una conferma istantanea
di riconoscibilità nel mondo fisico.

\subsection{Privacy}
La privacy è il prerequisito fondamentale per creare un
ambiente sicuro e confortevole per tutti gli utenti.
\paragraph{Direttiva UI/UX}
Il design deve enfatizzare attivamente l'assenza di
tracciamento GPS. L'interfaccia non deve mai presentare radar
o mappe che mostrino la posizione esatta degli utenti.
\paragraph{Obiettivo}
La tecnologia BLE deve essere percepita come una zona \textit{judgment-free},
assicurando costantemente all'utente che la sua posizione è
invisibile e che il controllo dell'interazione rimane totalmente
nelle sue mani.

\section{Archetipo Primario}
L'archetipo principale che definisce l'identità del brand è:
\[
    \textbf{The Everyman (L'Uomo Comune)}
\]
La sua missione è l'appartenenza. L'app è lo strumento quotidiano
e normale per connettersi con l'ambiente circostante, che
tu voglia salutare il tuo gruppo di amici al tavolo
accanto o rompere il ghiaccio con uno sconosciuto al
bancone Agisce come scudo di normalizzazione sociale, abbassando
le barriere all'ingresso e ponendo l'app come uno strumento
democratico di comunità. In questo modo, trasforma l'interazione
da evento eccezionale o potenzialmente imbarazzante in un gesto
semplice, leggero e accessibile, restituendo spontaneità all'incontro
e rendendo più facile partecipare a ciò che accade intorno.
L'app non forza la connessione, la rende soltanto possibile,
più fluida e più umana, creando un ecosistema in cui amici,
conoscenti e sconosciuti possono coesistere nello stesso
spazio fisico come parte di una medesima energia
collettiva. Così, l'app non si presenta come un mezzo per
"rimorchiare", ma come un'infrastruttura sociale discreta che
facilita la vicinanza, rafforza il senso di presenza ed
aiuta le persone a sentirsi parte di un momento condiviso.

\section{Archetipo Secondario}
L'archetipo secondario che definisce l'identità del brand
è:
\[
    \textbf{The Lover (L'Amante)}
\]
La sua missione è il soddisfare il bisogno profondo di validazione
emotiva, vicinanza umana ed esperienza sensoriale. Il suo ruolo
emerge nel momento in cui la connessione smette di essere soltanto
una possibilità sociale astratta e assume un valore più personale,
immediato e percepito. Se l'Everyman normalizza l'incontro
e lo rende accessibile, il Lover ne definisce invece la
qualità emotiva: trasforma la prossimità in presenza significativa,
la disponibilità in attenzione reciproca e l'interazione
in un'esperienza vissuta nel mondo reale. In questo modo,
l'app non si limita a facilitare il contatto tra persone
nello stesso spazio fisico, ma attribuisce a quel contatto
una densità relazionale maggiore, che può tradursi in
amicizia, networking, attrazione o semplice riconoscimento
umano. Il Lover interviene quindi come archetipo secondario
incaricato di dare profondità, intimità e memorabilità all'esperienza,
senza mai alterare il posizionamento pubblico del brand come
strumento inclusivo, quotidiano e orientato alla community.

\section{The Safe Space}
La tecnologia BLE e la politica di "Zero GPS / Zero Tracking"
rappresentano il punti di sintesi più chiaro tra i due
archetipi che strutturano l'identità del brand.

\paragraph{Dal lato dell'Everyman}
l'assenza di mappe e di tracciamento continuo definisce una
\textit{judgment-free zone}, riduce il rischio percepito di
esposizione e contribuisce a creare un amibente più neutro,
paritario e socialmente confortevole per tutti gli utenti.
\paragraph{Dal lato del Lover}
la privacy assoluta costruisce invece uno spazio protetto
di intimità relazionale, necessario affinchè le persone possano
abbassare le difese, mostrarsi più disponibili e compiere
il primo gesto di apertura sociale. In questo modello, la
tecnologia non si impone come elemento visibile dell'esperienza,
ma agisce come infrastruttura silenziosa che rende possibile
la connessione, lasciando emergere in primo piano soltanto
la presenza reciproca e la qualità del contatto umano.

% Keep everything below this line at the end of the document
% ===========================================================================
\section{Flusso \dots}
\dots

\section{Threat model e mitigazioni}

Questa sezione identifica i principali rischi strategici
connessi all'architettura archetipica del brand e le
relative mitigazioni. I threat model qui elencati non riguardano
vulnerabilità tecniche dell'applicazione, ma possibili
derive percettive, narrative o progettuali che potrebbero compromettere
la coerenza dell'identità definita nel presente documento.

\subsection{Percezione di dating app}

\textbf{Descrizione del rischio:} Un'eccessiva enfasi sull'archetipo
\textit{The Lover} può spingere il brand verso una lettura
prevalentemente romantica o seduttiva, facendo percepire Wavvy
come una dating app anziché come una piattaforma sociale di
prossimità orientata alla community.

\textbf{Impatto:} Riduzione del mercato indirizzabile, aumento
dello stigma sociale nell'utilizzo pubblico dell'app,
ambiguità nel posizionamento e perdita di credibilità come
infrastruttura relazionale generalista.

\textbf{Mitigazioni:}
\begin{itemize}
    \item Mantenere \textit{The Everyman} come archetipo
          primario esplicito in tutte le superfici pubbliche del
          brand.

    \item Evitare nella comunicazione e nel design riferimenti
          visivi o verbali troppo legati a seduzione, flirt o
          romanticismo.

    \item Descrivere sistematicamente l'app come strumento di
          community, prossimità sociale e connessione quotidiana.

    \item Validare ogni scelta di copy, UI e visual identity
          rispetto alla domanda: \enquote{Questa scelta fa sembrare Wavvy una dating app?}
\end{itemize}

\subsection{Percezione di utility fredda e
    priva di emozione}

\textbf{Descrizione del rischio:} Un'eccessiva enfasi sull'archetipo
\textit{The Everyman} può rendere il brand troppo neutro,
funzionale o descrittivo, facendo percepire Wavvy come una
semplice utility di rilevamento persone nelle vicinanze.

\textbf{Impatto:} Bassa memorabilità, debole coinvolgimento
emotivo, riduzione della ritenzione e impoverimento del
valore percepito dell'esperienza.

\textbf{Mitigazioni:}
\begin{itemize}
    \item Preservare il ruolo del \textit{Lover} come
          archetipo secondario incaricato di dare profondità
          emotiva all'interazione.

    \item Tradurre il concetto di prossimità in presenza significativa,
          non in semplice informazione funzionale.

    \item Favorire una narrazione del prodotto centrata su esperienza,
          riconoscimento umano e qualità del contatto.

    \item Verificare che il brand non venga ridotto a lessico
          esclusivamente tecnico, logistico o utility-driven.
\end{itemize}

\subsection{Ambiguità archetipica}

\textbf{Descrizione del rischio:} La compresenza di due
archetipi può generare un'identità confusa se il confine
tra ruolo primario e ruolo secondario non viene mantenuto
in modo disciplinato.

\textbf{Impatto:} Incoerenza tra prodotto, marketing e
design; difficoltà decisionali interne; perdita di
leggibilità del brand verso utenti, stakeholder e partner.

\textbf{Mitigazioni:}
\begin{itemize}
    \item Formalizzare che \textit{The Everyman} rappresenta
          il volto pubblico del brand, mentre \textit{The Lover}
          ne definisce la profondità esperienziale.

    \item Stabilire criteri di utilizzo dei due archetipi nei
          diversi contesti: comunicazione pubblica, onboarding,
          UX, meccaniche relazionali, visual design.

    \item Utilizzare il presente documento come riferimento vincolante
          per approvare scelte narrative e creative.

    \item Introdurre una regola di governance secondo cui ogni
          decisione di brand deve poter essere ricondotta chiaramente
          a uno dei due ruoli archetipici.
\end{itemize}

\subsection{Contraddizione tra promessa di
    privacy e percezione del prodotto}

\textbf{Descrizione del rischio:} Se la promessa di \enquote{Safe Space}
non viene percepita come credibile, l'intero impianto archetipico
perde stabilità, poiché sia \textit{The Everyman} sia \textit{The
    Lover} dipendono dalla riduzione del rischio sociale.

\textbf{Impatto:} Calo di fiducia, aumento della
resistenza all'adozione, riduzione della propensione a compiere
il primo gesto di apertura sociale e indebolimento del posizionamento
differenziante.

\textbf{Mitigazioni:}
\begin{itemize}
    \item Rendere la privacy un principio strutturale e visibile
          del prodotto, non una semplice promessa di marketing.

    \item Comunicare in modo chiaro l'assenza di GPS, mappe
          e tracking continuo.

    \item Assicurare che il tono visivo e verbale non evochi
          sorveglianza, esposizione o geolocalizzazione
          aggressiva.

    \item Verificare che tutte le interazioni chiave dell'app
          trasmettano controllo, discrezione e assenza di
          pressione.
\end{itemize}

\subsection{Traduzione visiva incoerente con
    il blend archetipico}

\textbf{Descrizione del rischio:} Una direzione cromatica o
visiva troppo spostata verso un solo polo del blend può alterare
la percezione del brand. Colori troppo aggressivi, troppo romantici
o troppo istituzionali possono compromettere l'equilibrio tra
accessibilità e intensità emotiva.

\textbf{Impatto:} Distorsione percettiva del
posizionamento, perdita di coerenza tra identità strategica
e interfaccia, associazione involontaria con dating app o
utility corporate.

\textbf{Mitigazioni:}
\begin{itemize}
    \item Utilizzare la direzione cromatica come conseguenza
          del blend, non come scelta autonoma o puramente
          estetica.

    \item Escludere famiglie cromatiche già identificate come
          incoerenti con il nucleo archetipico.

    \item Validare ogni proposta visuale rispetto ai due assi
          strategici: inclusione sociale e magnetismo
          relazionale.

    \item Mantenere una distinzione chiara tra colore principale,
          colori di supporto e accenti funzionali.
\end{itemize}

\subsection{Sbilanciamento sul contesto romantico
    a discapito della community}

\textbf{Descrizione del rischio:} L'archetipo \textit{The
    Lover}, se applicato in modo non controllato, può spingere
l'esperienza a privilegiare scenari di attrazione romantica
rispetto a scenari di amicizia, community e networking.

\textbf{Impatto:} Riduzione della pluralità d'uso dell'app,
restringimento del pubblico, perdita della promessa di piattaforma
inclusiva e quotidiana.

\textbf{Mitigazioni:}
\begin{itemize}
    \item Descrivere il valore relazionale di Wavvy in termini
          ampi: amicizia, community, networking, presenza,
          scoperta sociale.

    \item Evitare che il lessico del brand si concentri esclusivamente
          su desiderio, attrazione o scelta romantica.

    \item Garantire che il brand racconti sempre una gamma estesa
          di connessioni umane possibili.

    \item Mantenere la community come tela pubblica del brand
          e l'intimità come qualità secondaria dell'esperienza.
\end{itemize}

\subsection{Mancanza di trasferibilità interna
    del framework archetipico}

\textbf{Descrizione del rischio:} Se il modello archetipico
resta troppo astratto o teorico, i team di design, prodotto
e comunicazione potrebbero interpretarlo in modi divergenti,
riducendone il valore operativo.

\textbf{Impatto:} Disallineamento tra team, incoerenza
esecutiva, decisioni estetiche o narrative arbitrarie,
perdita di efficacia del framework.

\textbf{Mitigazioni:}
\begin{itemize}
    \item Tradurre il modello archetipico in principi operativi,
          criteri decisionali e regole di esclusione.

    \item Collegare esplicitamente ogni archetipo a implicazioni
          di prodotto, UX, visual design e messaging.

    \item Usare esempi concreti di \enquote{coerente} e
          \enquote{incoerente} nel brand book e nei documenti di
          handoff.

    \item Prevedere una fase di allineamento interfunzionale
          prima dell'implementazione esecutiva.
\end{itemize}

\section{Alternative considerate e modelli mitigati}

Questa sezione documenta il razionale strategico alla base
della scelta di un'architettura archetipica ibrida per il brand.
L'obiettivo è motivare, in termini di business, perché l'utilizzo
di un singolo archetipo risulterebbe subottimale rispetto
a una configurazione mista, più adatta a sostenere scalabilità,
acquisizione utente, ritenzione emotiva e posizionamento di
lungo periodo.

\subsection{\textit{The Lover} come archetipo unico}

L'adozione esclusiva dell'archetipo \textit{The Lover} porterebbe
il brand verso un posizionamento eccessivamente orientato al
dating, assimilabile alla formula \enquote{app per incontrare persone attraenti nelle vicinanze}.
Tale configurazione presenta tre criticità principali.

In primo luogo, determinerebbe una contrazione del \textit{Total
    Addressable Market (TAM)}. Il mercato delle applicazioni di
dating puro è già saturo e si rivolge prevalentemente a un
segmento specifico di utenti, ossia single in cerca di
incontri. Un prodotto che ambite a diventare una piattaforma
sociale di ampia adozione deve invece poter includere utenti
single, persone già in relazione, gruppi di amici,
colleghi, community locali e contesti ibridi di socialità.
Un posizionamento esclusivamente \textit{Lover} ridurrebbe
quindi in modo significativo l'ampiezza del mercato potenziale.

In secondo luogo, un'impostazione puramente \textit{Lover}
aumenterebbe l'attrito sociale legato all'utilizzo del prodotto.
Le applicazioni percepite come strumenti per \enquote{rimorchiare}
tendono a essere usate in modo nascosto o con imbarazzo,
soprattutto in contesti pubblici, professionali o familiari.
Poiché Wavvy si basa sulla prossimità fisica in tempo
reale, la percezione di ambiguità o di esposizione sociale
costituirebbe un ostacolo strutturale all'adozione e all'utilizzo
spontaneo dell'applicazione nello spazio pubblico.

In terzo luogo, un modello esclusivamente \textit{Lover} comprometterebbe
la capacità dell'app di fungere da infrastruttura relazionale
generalista. In una simile cornice, anche un gesto neutro come
inviare un \textit{Wave} a un conoscente, a un vecchio amico
o a una persona con cui si desidera semplicemente aprire un
contatto non romantico rischierebbe di essere interpretato
come ambiguo o inappropriato. Ciò limiterebbe fortemente
la funzione di unificazione sociale che il prodotto
intende svolgere.

\subsection{\textit{The Everyman} come archetipo unico}

L'adozione esclusiva dell'archetipo \textit{The Everyman}
porterebbe invece il brand verso un posizionamento
puramente funzionale, assimilabile a una \enquote{utility radar per vedere chi c'è nei dintorni}.
Anche questa configurazione presenta limiti rilevanti.

Il primo limite riguarda la ritenzione emotiva. Le
piattaforme costruite unicamente sull'archetipo \textit{Everyman}
tendono a essere percepite come strumenti utilitari, più
che come esperienze memorabili. In assenza di una componente
emotiva distintiva, l'interesse iniziale dell'utente
rischia di esaurirsi rapidamente, con conseguente aumento del
\textit{churn rate}. Un'applicazione sociale basata solo sull'utilità
informativa fatica a costruire un'abitudine d'uso ricorrente.

Il secondo limite è l'assenza di intensità emotiva. Un
social network scala non soltanto quando informa, ma quando
attiva curiosità, aspettativa, riconoscimento e gratificazione.
Un approccio \textit{Everyman} puro descrive la presenza di
persone nello spazio, ma non attribuisce a tale presenza
un valore relazionale sufficiente a generare
coinvolgimento, memoria o desiderio di ritorno.

Il terzo limite riguarda il rischio di sterilizzazione delle
interazioni. Se il prodotto si posiziona in modo
eccessivamente neutro, amichevole e comunitario, finisce
per neutralizzare proprio quella tensione umana che rende
significative molte connessioni nel mondo reale. In tale scenario,
l'app rischia di trasformarsi in una piattaforma corretta
ma piatta, incapace di abilitare relazioni più dense,
spontanee e memorabili, siano esse amicali, professionali
o romantiche.

\section{Sintesi}

L'identità di Wavvy si fonda su un'architettura archetipica ibrida
costruita attorno a due poli complementari: \textit{The Everyman}
come archetipo primario e \textit{The Lover} come archetipo
secondario. Il primo definisce il volto pubblico del brand,
posizionando l'app come uno strumento inclusivo, quotidiano e
socialmente accessibile; il secondo conferisce profondità
emotiva, intensità relazionale e memorabilità all'esperienza,
senza spostare il brand verso una percezione esclusivamente
romantica o seduttiva.

Questo equilibrio permette a Wavvy di presentarsi come
infrastruttura sociale di prossimità, capace di facilitare
connessioni umane reali in modo naturale, discreto e privo
di attrito. La promessa di \enquote{Safe Space}, resa credibile
dall'assenza di GPS, mappe e tracking continuo e dal ricorso
alla tecnologia BLE, costituisce il punto di sintesi più forte
tra i due archetipi: da un lato crea una zona neutra,
paritaria e \textit{judgment-free}; dall'altro rende possibile
un'apertura relazionale più autentica, protetta e significativa.

Dal punto di vista strategico, il framework scelto consente
di evitare due derive opposte ma ugualmente dannose: quella
di una dating app percepita come ambigua o stigmatizzante,
e quella di una utility fredda, funzionale e priva di valore
emotivo. Ne consegue che tutte le future decisioni di brand,
UX, visual design e comunicazione dovranno preservare questo
bilanciamento, mantenendo la community come cornice pubblica
dell'identità e l'intimità relazionale come qualità distintiva
dell'esperienza.

\end{document}